\section{SpaceXAI / xAI、Grok 4.7を公式リリース}
SpaceXAIはGrok 4.7を発表し、4.6と同一価格・速度での改善、困難タスクでの長時間作業と自己確認、強化したセーフガードを主張した。Cursor、Grok Build、APIで利用可能とし、Artificial AnalysisやArenaによる初期評価も窓内で共有された。
\dailyxsource{SpaceXAI (@SpaceXAI)}{https://x.com/SpaceXAI/status/2102069815225586149}

\section{Xiaomi、MiMo-V2.6 Pro / Flashを公開}
Xiaomi MiMoはomnimodalのMiMo-V2.6 ProとFlashを発表し、オープンウェイト、技術報告、RL環境と学習コードの公開を表明した。ProについてはArtificial AnalysisのIntelligence Index 46でオープンモデル首位とする初期評価や、低価格性が話題となった。
\dailyxsource{Xiaomi MiMo (@XiaomiMiMo)}{https://x.com/XiaomiMiMo/status/2102138559952290106}

\section{OpenAI、数学とAIの独立諮問グループを発表}
OpenAIは数学者による独立諮問グループと協働し、数学成果の評価・発信、学術基準、研究・学習ツールについて助言を受けると発表した。IASがホストし、グループは独立運営で未依頼の助言や公開コメントも可能と説明されている。
\dailyxsource{OpenAI (@OpenAI)}{https://x.com/OpenAI/status/2102093145051943229}

\section{Hugging Face事案をめぐる責任論が再燃}
米財務長官Scott Bessentが、Hugging Face事案ではエージェントではなく人間側が責任を負う趣旨を述べたとするクリップが拡散した。Andrew Ngも同日、サンドボックスや監視、ツール利用者・製作者の責任を論じ、一時停止論に反対する長文を投稿した。
\dailyxsource{NIK (@ns123abc)}{https://x.com/ns123abc/status/2102055361419108806}
\dailyxnote{CNBC原クリップの公式ページはGrok実行時点で未確認。}

\section{SGLang、Blackwell向けNVFP4 KV cacheを解説}
LMSYSはSGLang上の4-bit NVFP4 KV cacheを紹介し、Blackwell環境でKVフットプリントを減らしつつ長コンテキストdecodeを高速化できると説明した。Qwen3.5-397B-A17Bで精度をほぼ維持したとの主張も示された。
\dailyxsource{LMSYS Org (@lmsysorg)}{https://x.com/lmsysorg/status/2102076478183907705}
\dailyxnote{ブログ自体は観測窓より前の日付で、窓内では公式解説投稿が観測された。}

\section{Qwen-Image-2.1、デモと推論スタック対応を公式拡散}
Qwen公式はHugging AppsのSpacesデモとSGLang-Diffusion対応を引用し、単一チェックポイントでの生成・編集・RGBA対応を案内した。生成物の権利に関するライセンス解釈も補足された。
\dailyxsource{Qwen (@Alibaba\_Qwen)}{https://x.com/Alibaba_Qwen/status/2101866671900401913}
\dailyxnote{被引用のSGLang day-0投稿は観測窓外。}

\section{Artificial Analysis / Arena、新モデルの初期評価を公開}
同一観測窓でGrok 4.7とMiMo-V2.6-Proの初期スコアが連続公開され、価格対性能とエージェント性能が比較された。Arena側はAutoEvalであることを明記しており、人間投票の収束前の速報値として読む必要がある。
\dailyxsource{Artificial Analysis (@ArtificialAnlys)}{https://x.com/ArtificialAnlys/status/2102074898327932987}

\section{Andrew Ng、AI危険論の過熱を批判}
Andrew Ngは直近のAI脅威論を批判し、最大の変化はサイバー能力にあるとの見方を提示した。Hugging Face事案についてはサンドボックスや監視の不備、並列実行の文脈を重視し、責任はツールではなく人にあるとしてAI開発の一時停止に反対した。
\dailyxsource{Andrew Ng (@AndrewYNg)}{https://x.com/AndrewYNg/status/2102140576498065758}

\section{Yann LeCun、LLM研究停止要求説を否定}
Yann LeCunは、LLMの研究開発停止を求めたことはないと投稿した。Galactica、OPT-175B、Llama群への関与を挙げつつ、LLMは有用だが人間水準AIへの本線ではないという従来の立場を維持した。
\dailyxsource{Yann LeCun (@ylecun)}{https://x.com/ylecun/status/2101992797398082008}

\section{Procedural Graphs、エージェント手続きを編集可能グラフへ}
二次解説投稿は、LLMエージェントの実行履歴から手続きグラフを編集する「Procedural Graphs」研究を紹介した。複数モデル・ベンチでの改善やMultiChallengeでの修復率向上が主張されたが、Grok実行時点ではarXiv本文や著者公式Xは未確認だった。
\dailyxsource{Rohan Paul (@rohanpaul\_ai)}{https://x.com/rohanpaul_ai/status/2101973711612190756}

\section{本番LLM推論トレース公開への二次言及}
窓内投稿は、Harvard等による1年分・61.2億リクエストのLLM推論メタデータトレース公開を紹介した。原発表は観測窓外であり、Chutes / Bittensorとの関係や分散MoEへの含意は二次投稿側の解釈として扱う。
\dailyxsource{Sami Kassab (@Old\_Samster)}{https://x.com/Old_Samster/status/2102013811771220400}

\section{NVIDIA、Grok 4.7祝賀とEinride連携を投稿}
NVIDIA公式はGrok 4.7をコーディング・知識労働向けモデルとして祝賀し、自社計算基盤による支援を表明した。別投稿ではEinrideがNVIDIA Hyperion / DRIVE上で次世代ドライバーを構築する発表も紹介された。
\dailyxsource{NVIDIA (@nvidia)}{https://x.com/nvidia/status/2102081129234751722}
